%======================================================================
\section{Certificates for the model}\label{app:certificates}\label{app:tables}
%======================================================================

The arithmetic is that of \cite{PaperI}. \emph{Exact} computations use integers and rationals (FLINT \texttt{fmpz},
\texttt{fmpq} and their polynomial and matrix types \cite{FLINT}). \emph{Interval} computations use Arb ball arithmetic
\cite{Joh17}: a real number is represented by a ball $[m\pm r]$ that is guaranteed to contain it, and every operation returns
a ball containing all possible results. Both are accessed through python-flint (version $0.9.0$); constants such as $\pi$
are balls, so the nodes $2\pi n$ are exact. Linear systems are solved by Arb's ball solver, which either returns balls
containing the exact solution of every system in the input balls, and thereby certifies that the exact matrix is
invertible, or fails. A sign is accepted only if the ball excludes $0$. Positivity of polynomials is certified by
coefficient signs.

\begin{lemma}[Positive-coefficient certificates]\label{lem:posc}
\textup{(a)} \textup{(Descartes at $c$.)} If $p\in\RR[t]$ and every coefficient of $p(c+t)\in\RR[t]$ is non-negative with
positive constant term, then $p>0$ on $[c,\infty)$. If $p$ vanishes to exact order $m$ at $c$ and the coefficients of
$t^m,t^{m+1},\dots$ in $p(c+t)$ are positive, then $p>0$ on $(c,\infty)$.

\textup{(b)} \textup{(Region certificates.)} Let $\mathcal R\subseteq\bigcup_i\phi_i([0,\infty)^m)$, where the coordinates of
each $\phi_i$ are rational functions with a common denominator $w_i>0$ on $[0,\infty)^m$. Let $G$ be a polynomial of
total degree $d$. If each polynomial $w_i^d\cdot(G\circ\phi_i)$ has non-negative coefficients and a positive constant
term, then $G>0$ on $\mathcal R$.

\textup{(c)} A polynomial with non-negative coefficients in $v\in[0,\infty)^m$ is non-negative there, and positive at every
point at which some monomial with a positive coefficient is positive.

With ball coefficients, ``positive'' means that the ball lies in $(0,\infty)$.
\end{lemma}

Parts (a) and (c) are also in \cite{PaperI}.

\begin{proof}
All three statements follow from the fact that a polynomial with non-negative coefficients takes non-negative values
on the non-negative orthant, with the indicated strictness. In (b), $w_i^d\cdot(G\circ\phi_i)$ is a polynomial because
$\deg G=d$, and $w_i>0$ transfers its sign to $G\circ\phi_i$.
\end{proof}

In the applications the coefficients are exact integers (Propositions~\ref{thm:PT2}--\ref{thm:PT4}, \ref{thm:W-J2},
\ref{thm:W-J3}) or balls (the Descartes tests in Section~\ref{ssec:toy-tail} and Proposition~\ref{prop:base}); Taylor shifts of
ball polynomials are computed by Horner's scheme at a ball containing $c$, which encloses the exact shifted coefficients.

\subsection{The certificates for small \texorpdfstring{$J$}{J}}\label{app:proofs-PT}

This subsection describes the computations behind Propositions~\ref{thm:PT2}--\ref{thm:PT4}; the logic was given after
Proposition~\ref{thm:PT4}. Throughout, $c=\frac{21}{20}$, $V=\prod_{i<j}(y_i-y_j)$, and all arithmetic is exact.

\emph{The polynomials $G_i$.} Let $\mathfrak M(y)$ be the $(2J+1)\times(2J+1)$ matrix whose rows are the data rows
$(\tau_k(y_j))_{k\le2J}$, $(\tau_k'(y_j))_{k\le2J}$, $j\le J$, followed by the row $(\tau_k(y))_{k\le2J}$. Its determinant
$\mathfrak Q(y)$ is a combination of the $\tau_k$, hence admissible, and vanishes doubly at each $y_j$; its constant
term is $\pm D$, where $D$ is the determinant of the system in \hyp{N}$_{\yv}$. The verifier computes the minors by
Laplace expansion along the pairs of node rows, divides $\mathfrak Q$ by $A_{\yv}$ exactly, checks that every coefficient
$N_i$ of the quotient is divisible by $V^4$, and sets $G_i:=\pm N_i/V^4$, with the sign chosen so that $G_0>0$ at
$y^*$. It then checks the identities (i) and (ii) of the proof of Proposition~\ref{thm:PT4} symbolically. For $J=3$ and
$J=4$ the fingerprints of the $G_i$ are printed in the logs.

\emph{Slack coordinates and region covers.} Write $y_1=c(3+A)$ and $y_{k+1}=y_k+c\,B_k$ with $B_1=B$, $B_2=E$,
$B_3=F$. The hypotheses of the propositions become $A\ge0$, $B,E,F>0$ and $2A+B\ge8$, $3A+2B+E\ge24$,
$4A+3B+2E+F\ge48$ (as far as they apply). All parameters $s,t,a,b,e,f$ below range over $[0,\infty)$.
\begin{itemize}
\item $J=2$: piece I, $A=4+s$, $B=t$; piece II, $A=4s/(1+s)$, $B=8-2A+t$. Every point with $A\ge4$ lies in I, and
every point with $A<4$ lies in II (take $s=A/(4-A)$).
\item $J=3$: PA, $A=8s/(1+s)+a$, $B=12/(1+s)+b$, $E=e$, which covers $3A+2B\ge24$ with $B>0$ (given $(A,B)$, any
$s\in[\max(0,12/B-1),A/(8-A)]$ works, and this interval is non-empty exactly when $3A+2B\ge24$); PA$'$, $A=8+a$; and
PB1, PB2, in which $(A,B)$ runs over the triangles with vertices $(0,8),(4,0),(8,0)$ and $(0,8),(8,0),(0,12)$ in the
projective barycentric form $\lambda=(1,s,t)/(1+s+t)$ with apex $(0,8)$, and $E=24-3A-2B+e$. The edges with
$\lambda_1=0$ lie in $\{B=0\}$ or in $\{3A+2B=24\}\subset$ PA.
\item $J=4$: if $\ell:=4A+3B+2E\ge48$ (which implies $3A+2B+E\ge24$), the pieces Ia1 ($B=16+b$), Ia$'$ ($A=12+a$)
and Ia3 (the triangle $(12,0),(0,16),(12,16)$ in $(A,B)$, with apex $(12,0)$) cover $4A+3B\ge48$, and Ib1, Ib2 (the triangles
$(0,8),(4,0),(12,0)$ and $(0,8),(12,0),(0,16)$, both with apex $(0,8)$, and $E=(48-4A-3B)/2+e$) cover the rest, each with
$F=f$. The edge omitted from Ia3 is $\{B=16\}$, which lies in Ia1; the edge omitted from Ib1 lies in $\{B=0\}$, outside the
region; the edge omitted from Ib2 is the segment $\{4A+3B=48\}$, which is not part of the domain $4A+3B<48$ of Ib1, Ib2 and
lies in Ia3 (its endpoint $(0,16)$ in Ia1). If $\ell<48$, the pieces II1--II6 are the cones from the apex $(0,8,8)$ (which is
the point $y^*$) over the six triangles that triangulate the facets $B=0$, $E=0$ and $\ell=48$ of the polytope
$\mathcal P:=\{A,B,E\ge0,\ 2A+B\ge8,\ 3A+2B+E\ge24,\ \ell\le48\}$, with $F=48-\ell+f$. The apex lies on the three other facets
$A=0$, $2A+B=8$ and $3A+2B+E=24$, so $\mathcal P$ is the union of the cones from the apex over the facets $B=0$, $E=0$,
$\ell=48$ (a convex polytope is the union of the cones from one of its vertices over the facets that do not contain it).
The parametrisation omits the base face of each cone, which lies in $\{B=0\}\cup\{E=0\}$, outside the region, or in
$\{\ell=48\}$, covered by the first part.
\end{itemize}
\emph{What is machine-checked.} For every piece, \nolinkurl{toy/pt/verify_pt.py} checks that the piece lies in the region
(the line \texttt{inside} of its logs). The converse inclusion, that the pieces cover the region, is the hand argument
above, for $J=2$, for $J=3$, and for the part $\ell\ge48$ of $J=4$. For the part $\ell<48$ of $J=4$ the script
\nolinkurl{toy/pt/check_pt4_region.py} checks in exact arithmetic that the eight vertices of $\mathcal P$ and the six
triangles are as stated and that the six cones have total volume $1024/3$, the volume of $\mathcal P$ computed
independently; together with the convexity argument above this gives the cover.
In each piece the coordinates are rational functions of the parameters with a positive common denominator, and the
verifier checks that each homogenised $G_i\circ\phi$ has only positive coefficients and a positive constant term
(Lemma~\ref{lem:posc}(b)); the numbers of polynomials and their sizes are in Table~\ref{tab:app-toycert}. No subdivision
beyond these pieces is needed: on the whole region the numerators are coefficientwise positive in the slack variables.

\subsection{Design of the certificates for Theorem~\texorpdfstring{\ref{thm:PP80}}{PP80}}\label{app:PP80}

\emph{Chain route.} For a chain $Z$ of level $K$, the coefficients $q_m$ of $\omega_Z\pf$, for a monic $\pf$ of degree
$K$, are affine in the $K$ unknown lower coefficients of $\pf$. By Proposition~\ref{prop:umbral}(c), $\omega_Z\pf$ is
admissible if and only if $\sum_mq_m|s(m,2i-1)|=0$ for $i=1,\dots,K$; the unsigned Stirling numbers are exact
integers. This gives a $K\times K$ linear system for $\pf_Z$ (and a $(K+1)\times(K+1)$ system for $\pg_Z=\pf_{Z+0}$), whose
successful ball solution proves regularity. The script then checks $(\mathrm{I0})_K$, the increments
$\varrho_{k+1}-\varrho_k>0$, and $\pg^{(K)}(z_{K+1})>0$, at $4000$ bits; the worst relative radius over all levels is below
$10^{-785}$.

\emph{Direct route.} Here the interpolants are computed in the basis $\tau_0,\dots,\tau_K$ of $\Asp^{\mathrm{pol}}_K$,
with $\tau_k$ generated exactly by $\tau_k=(-1)^k(\Eul-y)^{2k}\,1$ (the conjugate $e^y\Eul e^{-y}=\Eul-y$); the cofactors are
obtained by exact division by $\omega_Z$ in ball arithmetic and compared with the expected degrees and leading
coefficients. The script checks $(\mathrm{I0})_K$, $(\mathrm{I1}')_K$ in the cross-multiplied form
$g_{k+1}f_k-g_kf_{k+1}>0$, and $(\mathrm{Y0})_K$ for every $K\le160$; and, for every $J\le80$ and without appeal to
Theorem~\ref{thm:A}, \hyp{N} for the system with $P(0)=1$, $p_{2J}>0$, strict $\TEL_J$, positivity of the half-step
cofactor, and the consistency relation $p_{2J}\prod_jy_j^2\,\tilde R_J(0)=1$ as a ball containment. The level range is
split into two runs ($K\le134$ and $135\le K\le160$). At
$K=160$ about $1150$ decimal digits are lost, so $6000$ bits leave a wide margin.

The two routes share no code beyond the python-flint library and agree on every level: they certify the same signs,
and their margins agree to the printed digits.

\subsection{Tables}\label{ssec:app-tables}

\begin{table}[ht]
\centering\small
\caption{Computer-assisted statements about the model. ``Exact'': integer or rational arithmetic; ``Arb'': ball
arithmetic at the stated precision. Scripts are in \texttt{toy/} (Appendix~\ref{app:anc}).}\label{tab:app-toycert}
\begin{tabular}{@{}>{\raggedright\arraybackslash}p{0.2\textwidth}>{\raggedright\arraybackslash}p{0.17\textwidth}>{\raggedright\arraybackslash}p{0.33\textwidth}>{\raggedright\arraybackslash}p{0.22\textwidth}@{}}
\toprule
Statement & Arithmetic & Size & Script \\
\midrule
Prop.~\ref{thm:PT2} & exact & $G_i$: $54,\dots,26$ terms; $10$ substituted polynomials; second route by Cramer's rule &
\nolinkurl{pt/verify_pt.py}, \nolinkurl{pt/verify_pt2_cramer.py} \\
Prop.~\ref{thm:PT3} & exact & $G_i$: $1304,\dots,636$ terms; $4$ pieces, $28$ polynomials, up to $38\,104$ terms &
\nolinkurl{pt/verify_pt.py} \\
Prop.~\ref{thm:PT4} & exact & $D$: $308\,386$ terms; $G_i$: $45\,612,\dots,22\,604$ terms; $11$ pieces, $99$ polynomials,
up to $219\,336$ terms; region cover checked separately & \nolinkurl{pt/verify_pt.py}, \nolinkurl{pt/check_pt4_region.py} \\
Prop.~\ref{prop:AP-B-cert} & exact & $K\le90$; up to $181$ coefficients & \nolinkurl{ap/verify_ap.py},
\nolinkurl{ap/verify_halfstep.py} \\
Thm.~\ref{thm:PP80} & Arb, $4000$ and $6000$ bits & levels $K\le160$, $J\le80$; two independent routes &
\nolinkurl{pp80/verify_pp80_chain.py}, \nolinkurl{pp80/verify_pp80_tau.py} \\
Prop.~\ref{prop:PP-dominates} & exact & $K\le2.17\cdot10^5$ & \nolinkurl{pp/verify_pp_prefix.py} \\
Prop.~\ref{prop:W-false} & Arb, $1500$--$7000$ bits & up to $J=50$ & \nolinkurl{w/verify_w_cex.py} \\
Props.~\ref{thm:W-J2}, \ref{thm:W-J3} & exact & two routes; $J=3$: up to about $37\,000$ terms & \nolinkurl{w/verify_w_small.py},
\nolinkurl{w/verify_w_insertion.py} \\
\bottomrule
\end{tabular}
\end{table}

\begin{table}[ht]
\centering\small
\caption{Margins in Theorem~\ref{thm:PP80} along the prime-power chain (direct route, $6000$ bits): certified lower bound
for the smallest relative increment $\min_k(\varrho_{k+1}-\varrho_k)/|\varrho_k|$ at level $K$, and certified upper bound for
the largest relative radius of a cofactor coefficient.}\label{tab:app-PP80}
\begin{tabular}{@{}rll@{}}
\toprule
$K$ & increment $\ge$ & radius $\le$ \\
\midrule
$20$ & $0.3869948$ & $10^{-1732}$ \\
$40$ & $0.2650057$ & $10^{-1609}$ \\
$60$ & $0.1690426$ & $10^{-1464}$ \\
$80$ & $0.1159777$ & $10^{-1311}$ \\
$100$ & $0.08632101$ & $10^{-1154}$ \\
$120$ & $0.06810364$ & $10^{-990}$ \\
$140$ & $0.05605644$ & $10^{-818}$ \\
$160$ & $0.04739370$ & $10^{-648}$ \\
\bottomrule
\end{tabular}
\end{table}

The certified increments decrease with $K$, roughly like $1/K$: the overall log--log slope between $K=20$ and $K=160$ is
about $1.01$ (local slopes $0.55$, $1.11$, $1.31$, $1.32$, $1.30$, $1.26$, $1.26$), and $K$ times the increment falls
monotonically from $10.6$ ($K=40$) to $7.58$ ($K=160$). Divided by the scale $K/(k^*(K-k^*))$ of Newton's inequalities at
the position $k^*\approx0.66K$ of the minimum, they fall from about $2.38$ ($K=40$) to $1.70$ ($K=160$) (computed from
the table). So, relative to the natural $1/K$ scale, the margin is still decreasing at $K=160$; nothing in the data
suggests a sign change, but nothing suggests a positive limit either.


\begin{table}[ht]
\centering\small
\caption{Counterexamples to W (Proposition~\ref{prop:W-false}), in the units $n=y/2\pi$; values rounded, signs certified.}
\label{tab:app-W}
\begin{tabular}{@{}>{\raggedright\arraybackslash}p{0.2\textwidth}>{\raggedright\arraybackslash}p{0.27\textwidth}>{\raggedright\arraybackslash}p{0.45\textwidth}@{}}
\toprule
Node set & Failure & Values \\
\midrule
$\{1,\dots,7\}$ & $\mu_7<0$ & $\mu_7\approx-2.332\cdot10^{-18}$ \\
$\{1,\dots,10\}$ & $\mu_9,\mu_{10}<0$; $p_1^{(10)}<p_1^{(9)}$ & $\mu_9\approx-8.924\cdot10^{-23}$, $\mu_{10}\approx-1.403\cdot10^{-25}$,
$p_1^{(10)}-p_1^{(9)}\approx-7.341\cdot10^{-11}$ \\
$\{1,\dots,5\}\cup\{m\}$ & $S<0$; $\mu_Y,\Delta<0$ & $S\approx-3.478\cdot10^{-6}$; $\Delta\approx-1.972\cdot10^{-9}$ ($m=10$),
$-1.773\cdot10^{-13}$ ($m=1000$) \\
$(1+\eps)j$, $j\le10$ & $\mu_9,\mu_{10}<0$ & $\eps=10^{-4}$: $S_{\le5}\approx-3.29\cdot10^{-6}$; $\eps=10^{-3}$: $S_{\le5}\approx-1.54\cdot10^{-6}$ \\
$(1+\frac1{200})j$, $j\le10$ & $\mu_{10}<0$ & $\mu_{10}\approx-3.025\cdot10^{-26}$ \\
$\frac{11}{10}+j$, $0\le j<50$ & $\mu_{48},\mu_{49},\mu_{50}<0$ & $\mu_{48}\approx-7.050\cdot10^{-122}$, $\mu_{50}\approx-1.720\cdot10^{-129}$ \\
\bottomrule
\end{tabular}
\end{table}
